\scanpage{441}
\mspage{223}

des autom. top. de $\widetilde U$ compatibles avec la
relation d'équivalence \unsure{définie} par $\widetilde U\unsure{\longrightarrow|-}U$, et
induisant sur $U$ l'\unsure{autom.} $\mathrm{id}_U$ ou $\sigma$.

\deleted{Pour} Alors

a) (Pour mémoire) $U^\sigma$ est une sous-variété fermée
de $U$ de dim 1.

b) Pour $\forall\ x\in U^\sigma$, on trouve une classe
de \inserted{$\pi$-}conjugaison de \deleted{éléments \unsure{de}}\inserted{scindages} de l'extension
$\mathcal E$, à la façon habituelle, d'où une
application
\deleted{(\unsure{4|5})}
\[
 \begin{aligned}
 U^\sigma\longrightarrow\mathrm{Sc}(\mathcal E,\mathbb Z/2)
 &=\text{ens. des classes de}\\
 &\quad\text{$\pi$-conj. de scindages de $\mathcal E$}\\
 &\quad\text{\unsure{sur} $\mathbb Z/2$, ou encore des}\\
 &\quad\text{éléments d'ordre $2$ de}\\
 &\quad\text{$\mathcal E$ qui ne sont pas dans $\pi$.}
 \end{aligned}
 \tag{5}
\]

Cette application se factorise en une
\uline{bijection}
\[
 \pi_0(U^\sigma)\simeq\mathrm{Sc}(\mathcal E,\mathbb Z/2)
 \tag{6}
\]

c) Soit $Z_i$ une comp. conn. de $U^\sigma$,
correspondant à un \deleted{scindage} $\sigma_i\in\mathcal E^-$ \deleted{tel}\inserted{d'}
ordre $2$. Alors \deleted{\illegible{1} $\pi$\illegible{1}}
\[
 \pi^{\sigma_i}=\{1\}\quad\text{\unsure{sss|ssi}}\quad Z_i\simeq\mathbb R,
 \quad\deleted{\unsure{\text{et}}}
 \tag{7}
\]

d) Si $Z_i\not\simeq\mathbb R$ i.e. \deleted{\illegible{4} $Z_i$} $Z_i\simeq\unsure{S^1|\mathbb S^1}$, alors,
pour une orientation choisie de $Z_i$, désig-

\scanpage{442}
nant par $g_i$ l'élément de $\pi=\pi_1(U)$
qu'il définit (défini à conj. près), et
par $\varphi_i:\mathbb Z\longrightarrow\pi$ l'hom $\varphi_i(n)=g_i^n$ de
$\mathbb Z$ dans $\pi$, on a :

$1^\circ)$ $\varphi_i$ est injectif

$2^\circ)$ quitte à remplacer $\varphi_i$ par un conjugué,
on a
\[
\unsure{(9)|(8)}\qquad \pi^{\sigma_i}=\operatorname{Im}\varphi_i
\]

\uline{Dém} On trouve l'application (5) en
\deleted{\unsure{prenant}}, pour $\forall\ x\in U^\sigma$, l'opération \inserted{\unsure{can.}} de
$\mathbb Z/2$ sur le rev. universel $\widetilde U(\unsure{x})$ ponctué
en $x$, et en \deleted{choisissant} prenant les
op. \unsure{correspondantes} sur $\widetilde U$, \unsure{déduites} par
les iso $\widetilde U\simeq\widetilde U(x)$. On voit que
les \deleted{\unsure{seules}} op. $\sigma'$\deleted{\illegible{1}} obtenues sur $\widetilde U$ \inserted{(pour $x$ fixé)}, sont
celles pour lesquelles \deleted{\illegible{1}} $\widetilde U^{\sigma'}$ a un
point au dessus de $x$ --- donc l'ens
des $x$ qui donnent naissance à \deleted{\unsure{une}}\inserted{la}
classe d'un $\sigma'$, sont les élts de
l'image de $\widetilde U^{\sigma'}$ par la projection
$\widetilde U\longrightarrow U$. Comme par le lemme $\widetilde U^{\sigma'}$
est non vide et connexe, il s'ensuit que

\scanpage{443}
\mspage{224}
son image dans $U^\sigma$ l'est aussi --- on va
voir que son image est exactement une comp.
connexe de $U^\sigma$ --- ce qui à la fois \unsure{prouvera}
la factorisabilité de (5) par \deleted{\unsure{le}} $\pi_0(U^\sigma)$
(\unsure{assez} triviale de ttes façons) et le fait que
l'application déduite (6) est \uline{bijective}.

Soit $\widetilde Z_i$ un \unsure{revêt\textsuperscript{t}|revêtement} universel de $Z_i$,
et prenons le rev. universel \unsure{correspondant}
$\widetilde U_i$ de $U$, de sorte qu'on a
\margin{\deleted{\diagram{\begin{tikzcd}[ampersand replacement=\&]
\widetilde Z_i \arrow[r] \arrow[d] \& \widetilde U_i \arrow[d] \\
Z_i \arrow[r] \& U
\end{tikzcd}}}}
\unsure{(8)|(9)}
\diagram{\begin{tikzcd}[ampersand replacement=\&]
\widetilde Z_i \arrow[r] \arrow[d] \& \widetilde U_i \arrow[d] \\
Z_i \arrow[r] \& U
\end{tikzcd}}
et par fonctorialité $\sigma$ opère aussi sur ce
diagramme (en opérant \unsure{trivial\textsuperscript{t}} sur $Z_i,\widetilde Z_i$),
soit $\sigma_i$ son opération sur $\widetilde U_i$. On
voit donc que $\widetilde Z_i$ s'envoie dans $\widetilde U_i^{\sigma_i}$,
qui s'envoie donc \uline{sur} $Z_i$ --- ce qui achève
déjà de prouver b).

\deleted{Supposons maintenant $Z_i$ \unsure{simple\textsuperscript{t}|simplement} connexe
i.e. $Z_i\simeq\mathbb R$, $\widetilde Z_i=Z_i$ --- donc $\widetilde Z_i\longrightarrow\widetilde U_i$ est
\unsure{une} section de $\widetilde U_i$.}
\inserted{Considérons \deleted{\unsure{l'image}} \inserted{l'image} \deleted{\unsure{induit}} de $\widetilde Z_i$ dans
$\unsure{\widetilde U_i|\widetilde U}\mid Z_i$, c'est une partie \uline{ouverte} et fermée}
(comme image d'un hom de \unsure{revêt\textsuperscript{ts}|revêtements}
étales sur la $\overline{\mathrm m}$ comp. $Z_i$), donc c'est

\scanpage{444}
\unsure{fortiori} une partie ouverte et fermée de
$\widetilde U_i^{\sigma_i}$, et comme cet espace est connexe,
il lui est égal.
\deleted{Lorsque}
\deleted{\[
(10)\qquad \widetilde U_i^{\sigma_i}=\operatorname{Im}(\widetilde Z_i\longrightarrow\widetilde U_i)
\]}
\inserted{De plus $\widetilde Z_i\longrightarrow\widetilde U^{\sigma_i}$ fait \deleted{\illegible{1}}\inserted{de}
$\widetilde Z_i$ un rev. étale de
son image $\widetilde U_i^{\sigma_i}$ dans $\widetilde U_i\mid Z_i$, $\ast$
et comme $\widetilde U_i^{\sigma_i}$ est
\unsure{simple\textsuperscript{t}|simplement} connexe,
on trouve \unsure{final\textsuperscript{t}|finalement}
\[
(10)\qquad \widetilde Z_i\overset{\sim}{\longrightarrow}\widetilde U_i^{\sigma_i}
\]}

Lorsque $Z_i$ est \unsure{simple\textsuperscript{t}|simplement} connexe i.e.
$\widetilde Z_i\simeq Z_i$, cela signifie aussi que
$\widetilde U_i^{\sigma_i}\longrightarrow Z_i$ est un homéomorphisme,
(donc $\widetilde Z_i=Z_i\longrightarrow\widetilde U_i^{\sigma_i}$ est l'\unsure{homéom.}
inverse). Comme, pour $x\in Z_i$,
\deleted{\unsure{les élts}} et $\widetilde x\in(\widetilde U_i^{\sigma_i})_x$, les
$\widetilde x'$ de $\widetilde U_i^{\sigma_i}$ au dessus de $x$ sont \deleted{\illegible{1}} \inserted{\unsure{les élém.}} de
la forme $\widetilde x\cdot\gamma$, avec $\gamma\in\pi_i^{\sigma_i}$
($\pi_i=\operatorname{Aut}(\widetilde U_i/U)$) il s'ensuit que l'on
a bien
\[
(11)\qquad \pi_i^{\sigma_i}=1
\]
ce qui est essentiellement la formule (7).

Dans le cas \deleted{général, \unsure{désignant}} \inserted{$Z_i$ non \unsure{simple\textsuperscript{t}|simplement} connexe,} posant
\[
(12)\qquad T_i=\operatorname{Aut}(\widetilde Z_i/Z_i)\underset{\unsure{f}}{\simeq}\pi_1(Z_i;\widetilde Z_i)
\]
on trouve un homomorphisme \unsure{can.}
\[
(13)\qquad \varphi_i:T_i\longrightarrow\pi_i
\]
et l'injectivité dans (10) équivaut au fait
que (13) est injectif. Bien entendu,

\scanpage{445}
\mspage{225}
le choix d'un générateur $g_i^{\circ}$ de $T_i$ équivaut au
choix d'une orientation de $Z_i$, et
$g_i=\varphi_i(g_i^{\circ})\in\pi_i$ est alors l'élément correspondant
de $\pi_1(U)$ dont il est question dans
l'énoncé (moins précis, \unsure{car} on n'y parle que
de classe de conjugaison). Il est clair
que $\sigma_i$ (\unsure{opérant} \unsure{trivial\textsuperscript{t}} sur $T_i$, et sur $\pi_i$
par transport de structure) commute à
l'hom injectif $\varphi_i$, donc \deleted{$\varphi_i(T_i)\subset$} $\varphi_i$
induit \deleted{$\varphi_i$} $T_i\longrightarrow\pi_i^{\sigma_i}$. Je dis que
c'est en fait un \uline{isomorphisme}
\[
\text{(14)}\qquad \varphi_i:T_i\xrightarrow{\sim}\pi_i^{\sigma_i}
\]
-- il reste à prouver la surjectivité. Mais
si $\gamma\in\pi_i^{\sigma_i}$, alors $\widetilde{x}\gamma\in\widetilde{U}_i^{\sigma_i}$, donc
$\widetilde{x}\gamma\in\varphi_i(\widetilde{Z}_i)$ ce qui signifie que
$\widetilde{x}\gamma=\widetilde{x}\varphi_i(g)$ \deleted{\unsure{donc}} pour $g\in T_i$ donc
$\gamma\in\operatorname{Im}\varphi_i$, cqfd.

\uline{Corollaire} 1 Soit $\sigma_i\in\mathcal{T}_g$ $(0\leq i\leq g)$,
correspondant à une anti-involution $\sigma_i$ de
$X_g$ telle que $\operatorname{card}\pi_0(X_g^{\sigma_i})=i$. Alors
\margin{et $X_g/\sigma$ non orientable....}
a) Il y a exactement $i$ classes de \inserted{$\Pi_g$}-conju-

\scanpage{446}
gaisons d'automorphismes effectifs d'ordre $2$
de $\Pi_g$, dans la classe $\sigma_i$ (ce qui prouve
que si $i\neq j$, $\sigma_i$ et $\sigma_j$ ne sont pas
conjugués dans $\mathcal{T}_g$ \deleted{\illegible{1}} \dots)

b) Si $u_i\in\mathcal{G}_g=\operatorname{Aut}_{\mathrm{lac}}(\Pi_g)$ est \inserted{d'ordre $2$,} dans la
classe $\sigma_i$, alors \deleted{$\unsure{\Pi_g}$}
\[
\text{(14)}\qquad \Pi_g^{\unsure{u_i}}\simeq\mathbb{Z}
\]

c) Si $u_i,u'_i\in\mathcal{G}_g$, \deleted{\unsure{satisfont} \unsure{tous deux} anti-}
\inserted{sont d'ordre $2$}
\deleted{involutions} de classe $\sigma_i$, alors $\exists\,\unsure{h}\in\mathcal{G}_g^+$
tel que
\[
\text{(15)}\qquad u'_i=hu_ih^{-1}=\operatorname{int}(h)u_i
\]
(\uline{NB} \unsure{Nécessairement}, l'image de $h$ dans $\mathcal{T}_g^+$
sera dans $(\mathcal{T}_g^+)^{\sigma_i}=\operatorname{Centr}_{\mathcal{T}_g}(\sigma_i)^+$.)

\uline{Démonstration} a) et b) sont des cas
particuliers du théorème, appliqué à
une anti-involution $\sigma_i$ de $X_g$, avec
$\pi_0(X_g^{\sigma_i})$ de cardinal $i$. Soit $Y_{\unsure{g,i}}=X_g^{\sigma_i}$,
surface compacte \inserted{à bord} non orientable \deleted{de} connexe
de type $(g-i,i)$, il est bien connu,
et immédiat, que le groupe des auto-

\scanpage{447}
\mspage{226}
morphismes \unsure{de l'|d'}\deleted{$\unsure{Y}$} une telle variété est transitif
sur $\pi_0(\partial Y)$, donc en \unsure{remontant} à $(X_g,\sigma_i)$, que
le groupe des autom de $(X_g,\sigma_i)$ est transitif
sur $\pi_0(X_g^{\sigma_i})$, \deleted{\illegible{2}} qu'on peut interpréter
aussi comme \unsure{l'ens} des classes de $\Pi_g$-conjugaisons
de $u_i$ comme dans b), c). Cela
implique donc qu'il existe $\dot{h}\in\unsure{\mathcal{T}_g^+|\mathcal{G}_g^+}$, \inserted{commutant à $\sigma_i$,} tel que
-- désignant par $h$ un \unsure{relevant|relèvement} dans $\mathcal{G}_g^+$ --
on ait $u'_i$ $\unsure{\Pi|\pi}$-conjugué à $\operatorname{int}(h).u_i$, ce
qui signifie aussi qu'on peut quitte à modifier
$h$), le choisir de façon qu'on ait (15)

\uline{Corollaire} 2 Sous les conditions du corollaire
précédent, \unsure{posant} $\Pi_g^{u_i}=T$ $(\simeq\mathbb{Z})$, on
a un diagramme de suites exactes
\[
\text{(16)}\qquad
\diagram{\unsure{
\begin{tikzcd}[ampersand replacement=\&,column sep=large]
1 \arrow[r] \& T \arrow[r] \arrow[d,equal] \&
\deleted{\unsure{\operatorname{Centr}_{\mathcal{G}_g}}}\mathcal{G}_g^{+u_i}
\arrow[r] \arrow[d] \&
\mathcal{T}_g^{+\sigma_i}
\arrow[d,hook] \quad \deleted{\unsure{\longrightarrow 1}}
\\
1 \arrow[r] \& T \arrow[r] \&
\mathcal{G}_g^{\deleted{+}u_i} \arrow[r] \&
\mathcal{T}_g^{\deleted{+}\sigma_i}
\quad \deleted{\unsure{\longrightarrow 1}}
\end{tikzcd}
}}
\]
et l'indice de $\mathcal{G}_g^{+u_i}$ dans $\mathcal{T}_g^{+\sigma_i}$
(et de $\mathcal{G}_g^{u_i}$ dans $\mathcal{T}_g^{\sigma_i}$) est $i$.
Comme les deux dernières flèches verticales
sont des inclusions de sous-groupes d'indice $2$,
il suffit de traiter l\deleted{\unsure{es}}'assertion concernant
$\mathcal{G}^+$, $\mathcal{T}^+$, \deleted{qui \illegible{2} \inserted{\illegible{1}} \illegible{1}}

\scanpage{448}
\deleted{1°) $\mathcal{G}_g^{+u_i}\longrightarrow\mathcal{T}_g^{+\sigma_i}$ est surjectif. Soit
$\unsure{\dot\gamma}\in\mathcal{T}_g^{+\sigma_i}$, alors $\unsure{u_i'}=\operatorname{int}(\gamma).u_i$ satisfait aux
conditions de c), donc $\exists h\in\mathcal{G}_g^{+\deleted{u_i}}$ tel que
\deleted{\illegible{1}} $u_i'=\operatorname{int}(h)u_i$ donc $\operatorname{int}(\unsure{h\gamma})u_i=u_i$ i.e.
$\unsure{h\gamma}\in$.}

Or $\mathcal{T}_g^{+\sigma_i}$ opère \inserted{transitivement} sur l'ens des
classes de $\Pi$-conjugaison des $u_i'$, d'après
c), \unsure{d'autre part} le stabilisateur \inserted{dans $\mathcal{T}_g^{+\sigma_i}$}, pour cette
action, de $u_i$, est formé des $\dot\gamma\in\mathcal{T}_g^{+\sigma_i}$
tels que \deleted{\unsure{il existe}} $\operatorname{int}(\gamma)u_i$ soit $\Pi$-conjugué
à $u_i$, i.e. tels qu'il existe $\alpha\in\Pi$, avec
$\operatorname{int}(\gamma)u_i=\operatorname{int}(\alpha)u_i$ i.e. $\alpha^{-1}\gamma\in\mathcal{G}_g^{+u_i}$, ce qui
signifie que $\dot\gamma$ est dans l'image de $\mathcal{G}_g^{+u_i}$,
cqfd.

\margin{\deleted{\unsure{comme}}}
\deleted{\uline{Corollaire 3} Soit $U$ une surface \deleted{compacte}
\inserted{\deleted{\illegible{1}} \unsure{0-conn.}, non \unsure{homéom.} au plan projectif réel.}
à bord, $\Pi=\pi_1(U)$, (Pour que \deleted{\illegible{1}} $\Pi$
ait des éléments de torsion, il \unsure{f et s.}
que $U$ soit non orientable et \deleted{$\partial U\ne\varnothing$. Alors}
\inserted{que $\partial U$ ait au \unsure{moins} une comp. connexe compacte.}
\deleted{\illegible{1}} Les éléments de torsion de
\deleted{\unsure{quand} $U$ \illegible{2} \unsure{orientable},}
$\Pi$ \deleted{\unsure{comme}} sont d'ordre 2, et leurs classes de
conjugaison sont en corr. 1-1 avec les
éléments de $\pi_0(\partial U)$ qui \unsure{correspondent} à des comp. \unsure{connexes} compactes.}

\scanpage{449}
\mspage{227}
J'ai envie de construire une figure \inserted{géométrique} où
on puisse mettre en évidence simultanément
des anti-involutions topologiques
qui \deleted{soient des} donnent naissance aux
$\sigma_i\in\unsure{\mathcal{T}_g^{\natural}|\mathcal{T}_g^{\sharp}}$ (à conj. près\unsure{..}), et aux divers
$u_i$ \inserted{$\in\mathcal{G}_g$} associés à un $\sigma_i$ (ce qui sera alors
facile, par la recette générale). Nous
savons que $\sigma_i$ \inserted{$\in\mathcal{T}_g$} s'obtient en regardant
un $X_g$ comme doublement orienté
d'un $Y=Y_{g-\deleted{\unsure{2}}i,i}$, \unsure{donc en partant}
du \deleted{\illegible{1}} plan projectif réel $Y_{\deleted{\unsure{g}}0}$, en
\margin{\deleted{\unsure{$D_0$}}}
y faisant $g$ trous, \deleted{dont} on \inserted{\deleted{\illegible{1}}} rebouche
$\unsure{g-i}$ par des rubans de Möbius, en
laissant les $i$ autres trous tels quels.
Soient $D_j$ ($1\le j\le g$) les disques disjoints
fermés \unsure{correspondants} aux "trous, donc
\deleted{\unsure{Alors}}
\[
(17)\qquad \boxed{V_{0,j}=Y_0\setminus\bigcup_j \overset{\circ}{D}_j}
\]
est une variété à bord
(non orientable), contenue à la fois dans
$Y_0$, et dans $Y=Y_{g-i,i}$, et coïncidant
$\overline{\mathrm m}$ avec $Y$ \unsure{en les points} de
\[
\boxed{\partial Y=\bigcup_{\unsure{g-i+1\le j\le g}}\partial D_j}.
\]
\deleted{Ainsi} Soit $X_0$ la sphère orientée qui \unsure{revêt} $Y_0$,

\scanpage{450}
et $U_0$ \deleted{\illegible{1}} $=X_0\vert V_0$, sa restriction sur $\bigcup D_j$,
c'est donc \inserted{le complémentaire d'} une réunion de $2g$ disques
\[
(18)\quad \deleted{\unsure{U_0}}\quad U_0\deleted{\unsure{U_g^i}}=X_0\setminus\bigcup_{1\le j\le g}\deleted{D}\Delta_j,
\]
où $\Delta_j\longrightarrow D_j$ est un rev. trivial à $2\deleted{\unsure{g}}$
feuillets de $D_j$ ($\Delta_j=\Delta'_j\amalg\Delta''_j\simeq D_j\times\unsure{\varepsilon_j|\mathcal{E}_j}$,
où $\unsure{\varepsilon_j|\mathcal{E}_j}$ est un ens à 2 éléments, qui s'identifie
à l'ens des deux orientations de
$D_j$\unsure{...}). \deleted{On \unsure{construit} \illegible{1} de sorte que}
\inserted{Soit \unsure{d'autre part}}
\deleted{\illegible{1}} $M_j$ \inserted{($1\le j\le g-i$)} le ruban de Möbius dont le bord \unsure{a|.}
été recollé à $V_0$ par un homéom.
\[
(19)\qquad \partial M_j\simeq\partial D_j
\]
\unsure{soit|sur}
\[
(20)\qquad M=\coprod_{1\le j\le g-i} M_i
\]
de sorte que
\[
(21)\qquad Y_{g\unsure{-|+}i,i}=V_0\amalg_{\partial M}M.
\]

\deleted{On voit qu'on a une appl. continue
\[
U_0\longrightarrow\unsure{X_{g-i,i}|X_{g,i}}
\]
qui induit un homéom.
\[
\operatorname{Int}(U_0)\longrightarrow Y_{g-i,i}\setminus(M\cup\partial Y)
\]
Plus précisément, l'image inverse}

Soit donc $X=X_g$ le \unsure{doublement|dédoublement} orienté
de $Y_{g-i,i}$ \unsure{défini} par (21), soit $\widetilde M=\coprod_{1\le j\le g-i}\widetilde M_j$
l'image inverse de $M$ dedans, \deleted{\unsure{$X_g$} \illegible{4}}
\inserted{un rev. étale de degré 2 de $M$,} \deleted{\illegible{3}}
\[
(\unsure{22|21})\qquad \deleted{X_g\unsure{\illegible{1}}}\quad X_{g,i}=X_g\setminus\operatorname{Int}(\widetilde M),
\]
